% !TEX program = xelatex
\documentclass[10pt,a4paper]{article}
\usepackage[margin=21mm]{geometry}
\usepackage{fontspec}
\setmainfont[Path=fonts/Main/,Extension=.otf,UprightFont=*-regular,BoldFont=*-bold,ItalicFont=*-italic,BoldItalicFont=*-bolditalic]{texgyretermes}
\usepackage{xcolor,enumitem,titlesec,tabularx}
\definecolor{accent}{HTML}{285F94}
\definecolor{muted}{HTML}{586A7D}
\usepackage[colorlinks=true,urlcolor=accent]{hyperref}
\setlength{\parindent}{0pt}
\setlength{\parskip}{3pt}
\setlist[itemize]{leftmargin=14pt,topsep=3pt,itemsep=2pt,parsep=0pt}
\titleformat{\section}{\large\color{accent}}{}{0pt}{}[\vspace{-3pt}\titlerule]
\titlespacing*{\section}{0pt}{13pt}{7pt}
\newcommand{\entry}[3]{\par\begin{tabularx}{\textwidth}{@{}Xr@{}}\textbf{#1}&{\small\color{muted}#2}\end{tabularx}\par{\small\color{muted}#3}\par}
\newcommand{\paper}[3]{\item #1.\ \textit{#2}.\ #3}
\hypersetup{pdftitle={Jiajun Wang - Curriculum Vitae},pdfauthor={Jiajun Wang}}
\pagestyle{plain}
\begin{document}
{\Huge Jiajun Wang}\hfill{\small\color{muted}Updated November 2026}\par
{\large PhD student in Mathematics}\par
Institut \'Elie Cartan de Lorraine (IECL), Universit\'e de Lorraine, Nancy, France\par
\href{mailto:jiajun.wang@univ-lorraine.fr}{jiajun.wang@univ-lorraine.fr}\enspace|\enspace\href{https://kazusa000.github.io}{kazusa000.github.io}\enspace|\enspace\href{https://github.com/kazusa000}{GitHub: kazusa000}

\section{Research interests}
Control theory, backstepping transformations, feedback design, numerical analysis, partial differential equations, optimization, and scientific computing.

\section{Education}
\entry{Universit\'e de Lorraine}{Oct. 2026 -- present}{PhD student in Mathematics, Institut \'Elie Cartan de Lorraine}
\entry{Universit\'e Paris-Saclay}{Sep. 2024 -- Sep. 2026}{Master studies in Mathematics and Applications: Analysis, Modeling, Simulation}
Coursework: PDE analysis, finite element and boundary element methods, finite- and infinite-dimensional optimization, and CPU/GPU high-performance computing.
\entry{Jiangnan University}{Sep. 2020 -- Jun. 2024}{BSc in Information and Computational Science; rank: 8/96}

\section{Research experience}
\entry{Backstepping transformations and feedback control}{Oct. 2026 -- present}{Doctoral research, IECL, Universit\'e de Lorraine}
\begin{itemize}
\item Study constructive transformations and feedback design for finite-dimensional linear systems with single or multiple inputs.
\item Prepare theoretical work as a manuscript and begin numerical comparisons, including invertible neural networks for computing transformations.
\end{itemize}
\entry{Shape optimization for hydrodynamic resistance in Stokes flow}{Apr. 2026 -- Aug. 2026}{Research internship, LS2N / CNRS, Nantes, France}
\begin{itemize}
\item Studied shape optimization methods for Stokes flow and hydrodynamic resistance.
\item Implemented numerical optimization using the Rodin finite element library.
\end{itemize}
\entry{Optimization methods for direction-of-arrival estimation}{Mar. 2024 -- Dec. 2024}{School of Science, Jiangnan University, China}
\begin{itemize}
\item Investigated numerical optimization for high-resolution signal estimation under mixed noise.
\item Contributed to model design and numerical evaluation.
\end{itemize}
\entry{Robust low-rank modeling and sparse noise removal}{Mar. 2023 -- May 2024}{School of Science, Jiangnan University, China}
\begin{itemize}
\item Studied iterative optimization for low-rank matrix recovery and signal decomposition.
\item Led numerical experiments and manuscript writing.
\end{itemize}
\entry{Parameter identification of a first-order RC equivalent circuit model}{Sep. 2022 -- Mar. 2023}{Undergraduate innovation research project, Jiangnan University, China}
\begin{itemize}
\item Designed an identification algorithm using numerical optimization and fractional-order gradient methods.
\item Performed data analysis, simulation validation, and manuscript writing.
\end{itemize}

\newpage
{\large Jiajun Wang}\hfill{\small\color{muted}Curriculum Vitae}\par
\section{Publications}
\textbf{Journal articles}
\begin{enumerate}[leftmargin=16pt,itemsep=10pt]
\paper{Xie, P., Chen, J., Cheng, L., \textbf{Wang, J.} (2026)}{Three-Variable Symmetric Alternating Direction Method of Multipliers for Non-Separable Convex Optimization}{Asian Journal of Control. \href{https://doi.org/10.1002/asjc.70174}{doi:10.1002/asjc.70174}.}
\paper{Xie, P., Chen, J., \textbf{Wang, J.}, Mao, Y., Zhu, G., Yi, Y. (2026)}{Alternating Direction Method of Multipliers for Direction of Arrival Estimation Under Mixed Noise}{Signal Processing, 247, 110648. \href{https://doi.org/10.1016/j.sigpro.2026.110648}{doi:10.1016/j.sigpro.2026.110648}.}
\paper{\textbf{Wang, J.}, Chen, J., Zhu, Q. (2024)}{The Application of the Accelerated Proximal Gradient Descent Algorithm for the Solution of the Weighted Schatten-$p$ Norm in Sparse Noise Extraction}{Circuits, Systems, and Signal Processing, 43, 5169--5190. \href{https://doi.org/10.1007/s00034-024-02697-z}{doi:10.1007/s00034-024-02697-z}.}
\end{enumerate}
\vspace{5pt}\textbf{Manuscript}
\begin{itemize}[itemsep=8pt]
\paper{Chen, J., \textbf{Wang, J.}, Gan, M., Xu, H., Na, J. (2026)}{Spectral Partitioning for Hierarchical Identification of Ill-Conditioned Systems: A Power-Iteration-Based Approach}{Manuscript.}
\end{itemize}
\vspace{5pt}\textbf{Conference paper}
\begin{itemize}[itemsep=8pt]
\paper{\textbf{Wang, J.}, Cheng, L., Chen, J. (2023)}{Parameter Identification of First-Order RC Equivalent Circuit Model Based on Improved Fractional-Order Gradient Descent}{35th Chinese Control and Decision Conference (CCDC), 293--298.}
\end{itemize}

\section{Software and personal projects}
\begin{itemize}[itemsep=6pt]
\item \textbf{Rodin and experiment interfaces.} Numerical research tools for finite element computation and Stokes shape optimization. \href{https://github.com/kazusa000/rodin}{Rodin}; \href{https://github.com/kazusa000/levelset-stokes-console}{levelset-stokes-console}.
\item \textbf{Innerlife.} Experimental virtual-person runtime with long-term memory, emotion, relationship state, and an inspection interface. \href{https://github.com/kazusa000/Innerlife}{GitHub repository}.
\end{itemize}

\section{Skills and languages}
\textbf{Programming:} Python, MATLAB, C++, C, Web development.\par
\textbf{Methods and tools:} finite elements, boundary elements, CPU/GPU computing, Lean.\par
\textbf{Languages:} Chinese (native), English (C1), French (B2), Japanese (basic).

\section{Other interests}
Violin (level 8/10), piano (beginner), and music arrangement. Self-taught arrangement and original music shared online.

\section{Contact}
\href{mailto:jiajun.wang@univ-lorraine.fr}{jiajun.wang@univ-lorraine.fr}\par
Additional academic contact: \href{mailto:jiajun.wang621@gmail.com}{jiajun.wang621@gmail.com}.
\end{document}
